%% IFT 6077 project proposal
%% ACM SIGCHI two-column format (acmart, sigconf). `nonacm` removes the
%% ACM reference block, copyright box and conference headers.
\documentclass[sigconf,nonacm]{acmart}

\usepackage{tikz}
\newcommand{\projectname}{TakeDo}
\usetikzlibrary{arrows.meta}

\begin{document}

\title{\projectname{}}
\subtitle{Giving Real-Time Production Notes to an Ensemble of AI Agents}

%% acmart stops with an error without \country; an empty one prints nothing.
\author{Travis Holt}
\affiliation{%
  \institution{Universit\'e de Montr\'eal}
  \city{}
  \country{}}
\email{travis.holt@umontreal.ca}

\begin{abstract}
Musicians think about their work in expressive, abstract, and relative terms (the verse should feel restless, the bass should sit under the kick). Yet Digital Audio Workstations (DAWs) only accept low-level inputs such as note velocities, timings and automation points spread across many tracks.
We propose \textbf{\projectname{}}, which reframes this work through the relationship between a \emph{record producer and a session team}: the producer gives notes while the track plays, and engineers and session players turn them into moves on the instruments and the console.
In \projectname{}, the user gives notes in speech or text during playback, and AI agents with production roles turn them into coordinated changes to MIDI notes, automation, tempo and structure. Direction is interpreted at three speeds: macro controls respond every frame through mappings the model selects, a fast model handles routine notes, and a large language model (LLM) plans complex ones. Every agent change passes through a change bus that gives the user priority and records it as editable project data, and the user's own edits serve as demonstrations the agents learn from.
More broadly, \projectname{} is a test case for an \emph{expressive interpretation layer}, which turns expressive direction into coordinated operations on a technical interface without giving up precision or editability.
\end{abstract}

\maketitle

%% ------------------------------------------------------------------
\section{Introduction}

A self-producing musician plays back a rough draft and hears that the chorus does not lift. The fix has no single control. It is spread across note velocities, timing, sounding parts, the balance between tracks, and the tempo. The musician has to work out which parameters to change, on which tracks, and by how much. However, DAWs only accept technical inputs. Realizing just one idea, such as ``let the bass breathe in the verse and push it in the chorus,'' means coordinated edits to many values that often interact. This is tedious for experts, out of reach for musicians without technical knowledge, and impossible in real time with two hands.

Current tools do not close this gap. Text-to-music generators such as Suno\footnote{\url{https://suno.com}} accept this kind of expressive description. However, they answer it by generating new audio as a closed artifact, instead of shaping an existing project in ways the musician can inspect and edit. LLM assistants can execute DAW edits~\cite{elkins2025dawzy}, but each request becomes an isolated edit made between playbacks.

In commercial studios, people close this gap. A record producer listens to playback and gives notes (``push the vocal into the chorus,'' ``play the hats looser''), and engineers and session players turn them into moves on the console and the instruments, often while the track plays. We propose \textbf{\projectname{}}, built on this \emph{producer and session team} metaphor. The user produces their draft as it plays, and an ensemble of AI agents plays the roles of the team. \projectname{} has three parts (Figure~\ref{fig:layers}).
\begin{itemize}
  \item \textbf{Input.} Like a producer on the talkback mic, the user gives notes by speech or text while the track loops, addressing a role, a track or a section (``drums, looser in the verse'' or ``make the chorus lift''). This happens in a \emph{producer view}, which also shows what each agent is doing. An interpretation layer turns each note into structured commands for the agents. The user can also ride \emph{macro controls}, faders whose meaning the model assigns (e.g., one fader becomes ``chorus intensity''). Other modalities, such as gesture, could later feed into the same layer.
  \item \textbf{Agent ensemble.} Agents take production roles. A performance agent shapes the velocity, timing and length of existing notes. An arrangement agent adds, removes or rewrites notes (fills, arpeggios, part thinning). A mix agent writes volume, pan, send, and effect automation. A conductor agent owns the tempo map and which parts play in each section. Changes, including rewritten notes, are applied ahead of the playhead, so the user hears them on the current pass. ``Bass under the kick'' becomes coordinated moves by several agents.
  \item \textbf{Change bus and performance capture.} Agents propose changes to a change bus. The bus applies them to the project state of openDAW\footnote{\url{https://github.com/andremichelle/openDAW}}, an open-source web DAW, and records each take as project edits. As in a DAW's \emph{touch} automation mode, when the user grabs a control in the timeline view, that parameter becomes \emph{human-held}. Agent writes to it are dropped until the user lets go, and the agent adapts.
\end{itemize}

\begin{figure*}[t]
  \centering
  %% System layers figure (Figure 1). \input inside a figure* environment.
%% An illustration, not an architecture diagram: box names and arrows only.
%% Orange = the user (input and views), blue = AI, gray = infrastructure.
\definecolor{hum}{RGB}{196,112,30}
\definecolor{humfill}{RGB}{253,242,229}
\definecolor{ai}{RGB}{42,106,155}
\definecolor{aifill}{RGB}{231,240,248}
\definecolor{infra}{RGB}{95,95,95}
\definecolor{bandfill}{RGB}{246,246,246}
\definecolor{bandline}{RGB}{214,214,214}
\definecolor{subtext}{RGB}{105,105,105}
\begin{tikzpicture}[
    x=1cm, y=1cm,
    font=\sffamily\scriptsize\bfseries,
    box/.style={draw, rounded corners=2pt, minimum height=0.8cm, align=center,
      inner sep=2pt, line width=0.6pt},
    human/.style={box, draw=hum, fill=humfill},
    ai/.style={box, draw=ai, fill=aifill},
    infra/.style={box, draw=infra, fill=white},
    arr/.style={-{Stealth[length=4pt]}, line width=0.6pt},
  ]

  %% Layer bands, with labels in the left gutter
  \foreach \y/\name in {0/Producer input, -1.55/Interpretation, -3.1/Agent ensemble,
                        -4.6/Change bus, -6.1/openDAW} {
    \draw[draw=bandline, fill=bandfill, rounded corners=3pt]
      (0,\y+0.55) rectangle (13.35,\y-0.55);
    \node[anchor=west, text width=1.9cm, font=\sffamily\tiny\bfseries, text=subtext]
      at (0.12,\y) {\MakeUppercase{\name}};
  }

  %% Producer input
  \node[human, minimum width=2.6cm] (macro) at (4.6,0) {Macro controls};
  \node[human, minimum width=3.6cm] (speech) at (8.45,0) {Speech or text};

  %% Interpretation
  \node[ai, minimum width=2.4cm] (map) at (4.6,-1.55) {Macro mappings};
  \node[ai, minimum width=2.3cm] (jev) at (8.45,-1.55) {Fast model};
  \node[ai, minimum width=2.1cm] (llm) at (12.1,-1.55) {LLM planner};

  %% Agent ensemble
  \foreach \x/\name in {3.6/Performance, 6.35/Arrangement, 9.1/Mix,
                        11.85/Conductor} {
    \node[ai, minimum width=2.6cm] at (\x,-3.1) {\name};
  }

  %% Change bus: a heavy bar, named by its layer label
  \node[infra, draw=black, line width=1.6pt, minimum width=10.85cm, minimum height=0.45cm]
    (bus) at (7.725,-4.6) {};

  %% openDAW
  \node[infra, minimum width=3.2cm] (state) at (3.9,-6.1) {Project state};
  \node[infra, minimum width=3.6cm] (takes) at (7.9,-6.1) {Captured takes};

  %% Views for the user: the producer view spans the input to ensemble layers,
  %% the timeline view spans the change bus and openDAW layers
  \node[human, minimum width=3.2cm, minimum height=4.2cm] (cview) at (15.4,-1.55)
    {Producer view};
  \node[human, minimum width=3.2cm, minimum height=2.6cm] (tview) at (15.4,-5.35)
    {DAW timeline view};

  %% Input -> interpretation
  \draw[arr, hum] (macro) -- (map);
  \draw[arr, hum] (speech) -- (jev);
  \draw[arr, ai] (jev) -- (map);
  \draw[arr, ai] (jev) -- (llm);

  %% Interpretation -> ensemble
  \draw[arr, ai] (map.south) -- (4.6,-2.55);
  \draw[arr, ai] (jev.south) -- (8.45,-2.55);
  \draw[arr, ai] (llm.south) -- (12.1,-2.55);

  %% Ensemble <-> change bus
  \draw[arr, ai] (4.0,-3.65) -- (bus.north -| 4.0,0);
  \draw[arr] (bus.north -| 9.4,0) -- (9.4,-3.65);

  %% Change bus <-> openDAW
  \draw[arr] (bus.south -| 3.2,0) -- (3.2,-5.7);
  \draw[arr] (4.6,-5.7) -- (bus.south -| 4.6,0);
  \draw[arr] (bus.south -| 7.9,0) -- (7.9,-5.7);

  %% Views
  \draw[arr] (13.35,0) -- (cview.west |- 0,0);
  \draw[arr] (13.35,-3.1) -- (cview.west |- 0,-3.1);
  \draw[arr, hum] (tview.west |- 0,-6.1) -- (13.35,-6.1);
\end{tikzpicture}

  \caption{(For illustration only. Not an architectural reference) System layers. Orange: the user; blue: AI components. Every change, from an agent or the user, passes through the change bus.}
  \Description{Layered diagram. Speech or text feeds a fast model, which directs an ensemble of performance, arrangement, mix and conductor agents, and passes complex requests to an LLM planner. Macro controls feed macro mappings, which the fast model selects and which drive the agents every frame. Agents exchange changes with a change bus, which updates the openDAW project state and captured takes. A producer view receives the user's input and the agents' activity, and a DAW timeline view sends the user's direct edits to openDAW.}
  \label{fig:layers}
\end{figure*}

\projectname{} draws on three theories seen in class. As an \emph{interface metaphor}~\cite{neale1997metaphors}, the producer and session team is a familiar model for directing many performers and operators at once. Giving notes in musical language keeps the \emph{semantic distance} small, because the user states the intent and not the parameters. It also keeps the \emph{articulatory distance} small, because notes take the form producers already use with engineers and players~\cite{hutchins1985direct}. Delegating to agents has long been seen as in tension with \emph{direct manipulation}. \projectname{} keeps both. Performance capture and the timeline view make every agent action visible, incremental and reversible~\cite{shneiderman1983direct}. Additionally the user can take over any control during playback, overriding the agents. Following \emph{mixed-initiative} principles~\cite{horvitz1999principles}, confident interpretations are carried out immediately, and uncertain ones wait for confirmation from subsequent input.

The core of \projectname{} is not specific to music. Video editors, lighting consoles and 3D animation tools also expose hundreds of technical parameters to users who think in expressive terms, and in those fields an art director plays the producer's role. We see \projectname{} as a test case for an \emph{expressive interpretation layer}. This is an AI layer that sits between a person's expressive direction and a technical interface. It turns that direction into coordinated low-level operations and leaves every operation visible and editable. We aim to make four contributions.
\begin{enumerate}
  \item \textbf{An architecture for AI interpretation of expressive direction at three speeds.} A fast structured-output model (``system one'') interprets routine notes. It chooses the roles and tracks addressed, an expressive descriptor and an intensity. Agents then turn descriptors into parameter recipes in code (e.g., \emph{looser} $\rightarrow$ timing offsets and velocity variation). Compound or novel notes, like ``have the bass answer the lead,'' go to an LLM. Its plans are validated and saved for reuse. Below both models, macro controls are applied every frame through mappings the models select, since riding a fader should have no detectable latency.
  \item \textbf{A mechanism for shared human--agent control of one interface.} The change bus gives the human priority at all times.
  \item \textbf{Direct edits as demonstrations of intent.} When the user edits the DAW directly, the agents observe the change through the bus. They use it to refine their understanding of what the user wants, so direct manipulation also teaches the agents, as in demonstrational interfaces~\cite{myers1992demonstrational}.
  \item \textbf{Delegation without loss of direct manipulation.} Agent actions are logged and become ordinary, editable artifacts of the underlying tool. The user can revise anything the agents do.
\end{enumerate}

\textbf{Our research question is whether expressive direction given during playback, interpreted at three speeds, lets musicians shape a draft with less effort than requests made between playbacks, without reducing their sense of control.} We also ask whether treating direct edits as demonstrations reduces how often later notes must be corrected. We will explore both with a working prototype on a short multitrack draft, comparing live direction with a turn-based version of the same system, and the ensemble with and without learning from edits.

%% ------------------------------------------------------------------
\section{Related Work}

\subsection{AI Assistance in Music Production}

Semantic audio tools link words such as \emph{warm} to effect settings, and automatic mixing systems set levels across tracks, but both make one kind of adjustment, outside of playback. DAWZY turns requests into reversible DAW edits~\cite{elkins2025dawzy}, each an isolated change made between playbacks. The Conductor Program~\cite{mathews1991radio} lets performers shape the tempo and dynamics of a stored score as it plays, but through fixed mappings. In \projectname{}, the user gives notes while the music plays, agents apply an intent across a passage and coordinate it across tracks, and the model chooses what each macro fader controls.

\subsection{Combining Language and Direct Manipulation}

DirectGPT applies direct manipulation principles to interaction with LLMs~\cite{masson2024directgpt}. Stylette turns natural-language styling goals into changes to many low-level CSS properties~\cite{kim2022stylette}. AI-Instruments reify prompts as reusable instruments that act on interface objects~\cite{riche2025aiinstruments}. \projectname{} brings these ideas to a time-based medium, where direction must be interpreted while the work plays and several agents act on the same project.

%% ------------------------------------------------------------------
\bibliographystyle{ACM-Reference-Format}
\bibliography{references}

\end{document}
